\section{The state space and the first anomaly}

A degree-one operator can define cohomology only when its degree-two square vanishes on the specified state space.
For the Virasoro algebra, this square is a quadratic ghost operator.
Its coefficient comes from a finite contraction in the ghost representation.
The zero modes then determine which restrictions form cochain complexes.

All vector spaces and tensor products below are algebraic over $\mathbb C$.
For homogeneous operators, the bracket is
$[A,B]_{\mathrm s}=AB-(-1)^{|A||B|}BA$.
Braces denote the bracket of two odd operators.
Matter operators have ghost degree zero.

Let $M$ be a Virasoro module with central charge $C\in\mathbb C$:
\begin{equation}\label{eq:vir-matter-relation}
 [L_m^M,L_n^M]=(m-n)L_{m+n}^M+
 \frac{C}{12}(m^3-m)\delta_{m+n,0}\,1.
\end{equation}
Assume that $M$ is restricted: for each $v\in M$, $L_n^Mv=0$ for all sufficiently large positive $n$.
No inner product or finite-dimensional weight-space hypothesis is imposed.

The ghost Clifford algebra has odd generators $b_n,c_n$, $n\in\mathbb Z$, and relations
\[
 \{b_m,c_n\}=\delta_{m+n,0},\qquad
 \{b_m,b_n\}=\{c_m,c_n\}=0.
\]
Its conformal vacuum $\Omega$ obeys
\begin{equation}\label{eq:vir-ghost-vacuum}
 b_n\Omega=0\quad(n\geq-1),\qquad
 c_n\Omega=0\quad(n\geq2).
\end{equation}
The ghost module is
\[
 \mathcal F=\Lambda\bigl(b_n\ (n\leq-2),\ c_n\ (n\leq1)\bigr)\Omega.
\]
A displayed generator acts by exterior multiplication.
Its Clifford partner acts by left exterior differentiation.
Thus the formulas specify the signs on every monomial.
Set $\deg\Omega=0$, $\deg c_n=1$, and $\deg b_n=-1$.
Set $E\Omega=0$ and assign both $b_n$ and $c_n$ energy $-n$.
The energy of a monomial is the sum of these assignments.

Normal ordering, denoted by colons, moves the creators in~\eqref{eq:vir-ghost-vacuum} to the left.
Each interchange of odd factors contributes a minus sign.
Contraction terms are omitted.
This convention applies to the whole monomial inside the colons.
Define
\begin{align}
 L_m^{\mathrm{gh}}&=\sum_{s\in\mathbb Z}(m+s):b_{m-s}c_s:,
 \label{eq:vir-ghost-L}\\
 R&=-\frac12\sum_{m,n\in\mathbb Z}(m-n):c_{-m}c_{-n}b_{m+n}:,
 \label{eq:vir-ghost-R}\\
 Q&=\sum_{n\in\mathbb Z}c_{-n}L_n^M+R
 \quad\hbox{on }\mathcal D=M\otimes\mathcal F.
 \label{eq:vir-Q}
\end{align}
Tensor factors commute because the matter action is even.
The degree of $Q$ is one.

\begin{lemma}[Algebraic domain]\label{lem:vir-domain}
Every sum in~\eqref{eq:vir-ghost-L}--\eqref{eq:vir-Q} has finitely many nonzero terms on each element of $\mathcal D$.
The same holds for
\[
 \nu_r=\frac12\sum_{a+b=r}(a-b)c_ac_b,
 \qquad
 K=\frac{C-26}{12}\sum_{n>0}(n^3-n)c_{-n}c_n.
\]
These operators and all their finite compositions preserve $\mathcal D$.
\end{lemma}

\begin{proof}
It suffices to use one matter vector and one ghost monomial.
For the matter term, indices $n\leq-2$ require the occupied creator $b_n$.
Only finitely many such indices occur.
The restrictedness of $M$ bounds the remaining positive indices.

For a normally ordered ghost monomial, every annihilator must remove an occupied creator.
Its indices therefore lie in a fixed finite set.
All creator indices have an upper bound: $-2$ for $b$ and $1$ for $c$.
The sum of the indices is fixed in each displayed operator.
Once the annihilator indices are fixed, this sum and the upper bounds give lower bounds for every creator index.
There are at most three factors.
Hence only finitely many terms act nontrivially.
The argument includes terms containing only creators.
It also applies to $\nu_r$ and $K$.
Each output is a finite linear combination of algebraic states, so compositions are defined on the same domain.
\end{proof}

For trivial matter, all $L_n^M$ vanish and $C=0$.
The first nonzero anomalous mode is visible without an infinite summation:
\begin{align*}
 R\Omega&=0, &
 Rb_{-2}\Omega&=-2b_{-2}c_0\Omega-b_{-3}c_1\Omega,\\
 R(b_{-2}c_0\Omega)&=4c_{-2}\Omega+b_{-3}c_0c_1\Omega
       +2b_{-2}c_{-1}c_1\Omega,\\
 R(b_{-3}c_1\Omega)&=5c_{-2}\Omega-2b_{-3}c_0c_1\Omega
       -4b_{-2}c_{-1}c_1\Omega.
\end{align*}
Consequently,
\begin{equation}\label{eq:vir-first-square}
 R^2b_{-2}\Omega=-13c_{-2}\Omega.
\end{equation}
The cancellation of the cubic output leaves a degree-two operator acting on a degree-minus-one input.
The calculation below determines every mode of this operator.

\section{The ghost central term}

\begin{lemma}[Ghost commutants]\label{lem:vir-commutant}
An operator on $M\otimes\mathcal F$ that supercommutes with every $b_n,c_n$ is zero if its ghost degree is nonzero.
An even operator on $\mathcal F$ commuting with every $b_n,c_n$ is a scalar.
\end{lemma}

\begin{proof}
The common kernel of the vacuum annihilators on $M\otimes\mathcal F$ is $M\otimes\mathbb C\Omega$.
Indeed, these annihilators are all the exterior derivatives of the creator variables.
A finite exterior polynomial killed by all these derivatives is constant.
A supercommuting operator sends this common kernel into itself.
The kernel has ghost degree zero, so an operator of nonzero ghost degree vanishes there.
Supercommutation with the creators makes it vanish on every state.
For an even operator on $\mathcal F$, its value on $\Omega$ is a scalar multiple of $\Omega$.
Commutation with the creators gives the same scalar on all monomials.
\end{proof}

\begin{proposition}[The ghost Virasoro algebra]\label{prop:vir-ghost-central}
On $\mathcal F$,
\begin{align}
 [L_m^{\mathrm{gh}},b_n]&=(m-n)b_{m+n}, &
 [L_m^{\mathrm{gh}},c_n]&=(-2m-n)c_{m+n},
 \label{eq:vir-ghost-transform}\\
 L_0^{\mathrm{gh}}&=E, &
 L_m^{\mathrm{gh}}\Omega&=0\quad(m\geq-1),
 \label{eq:vir-ghost-energy}\\
 [L_m^{\mathrm{gh}},L_n^{\mathrm{gh}}]
 &=(m-n)L_{m+n}^{\mathrm{gh}}
 -\frac{26}{12}(m^3-m)\delta_{m+n,0}\,1.
 \label{eq:vir-ghost-central}
\end{align}
\end{proposition}

\begin{proof}
The elementary contractions are
\[
 [:b_rc_s:,b_n]=\delta_{s+n,0}b_r,
 \qquad
 [:b_rc_s:,c_n]=-\delta_{r+n,0}c_s.
\]
They give~\eqref{eq:vir-ghost-transform}.
A term acting on $\Omega$ needs both factors to be creators.
For $m\geq0$ this is impossible.
For $m=-1$ the only possible term has coefficient zero.
The commutators of $L_0^{\mathrm{gh}}$ with the creators equal those of $E$.
Both operators kill $\Omega$, which proves~\eqref{eq:vir-ghost-energy}.

Let
$A_{m,n}=[L_m^{\mathrm{gh}},L_n^{\mathrm{gh}}]-(m-n)L_{m+n}^{\mathrm{gh}}$.
Substitution of~\eqref{eq:vir-ghost-transform} into the Jacobi identity shows that $A_{m,n}$ commutes with every ghost generator.
For example, its commutator with $b_r$ has coefficient
\[
 (n-r)(m-n-r)-(m-r)(n-m-r)-(m-n)(m+n-r)=0.
\]
For $c_r$, the coefficient is
\[
 (-2n-r)(-2m-n-r)-(-2m-r)(-2n-m-r)
 -(m-n)(-2m-2n-r)=0.
\]
Lemma~\ref{lem:vir-commutant} makes $A_{m,n}$ a scalar.
Its energy is $-m-n$, so this scalar vanishes unless $m+n=0$.

For $m\geq2$, the complete vacuum expression is
\[
 L_{-m}^{\mathrm{gh}}\Omega
 =\sum_{s=2-m}^{1}(s-m)b_{-m-s}c_s\Omega.
\]
Commuting $L_m^{\mathrm{gh}}$ through the first factor gives $(2m+s)b_{-s}$.
This contracts with $c_s$ to give $1$.
Its commutator with the second factor annihilates $\Omega$.
Therefore
\begin{align*}
 A_{m,-m}\Omega
 &=\sum_{s=2-m}^{1}(s-m)(2m+s)\Omega\\
 &=\left(\sum_{s=2-m}^{1}s^2
      +m\sum_{s=2-m}^{1}s-2m^2\cdot m\right)\Omega\\
 &=-\frac{13}{6}(m^3-m)\Omega.
\end{align*}
Here the interval has $m$ integers,
$\sum s=m(3-m)/2$, and
$\sum s^2=m(2m^2-9m+13)/6$.
The cases $m=0,1$ give zero by~\eqref{eq:vir-ghost-energy}.
Antisymmetry gives negative $m$.
This proves~\eqref{eq:vir-ghost-central}.
\end{proof}

Set $\mathcal L_n=L_n^M+L_n^{\mathrm{gh}}$ on $\mathcal D$.
The two summands commute.
Thus~\eqref{eq:vir-matter-relation} and~\eqref{eq:vir-ghost-central} give
\begin{equation}\label{eq:vir-total-central}
 [\mathcal L_m,\mathcal L_n]
 =(m-n)\mathcal L_{m+n}+A(m)\delta_{m+n,0},
 \qquad A(m)=\frac{C-26}{12}(m^3-m).
\end{equation}
This central term determines the square of $Q$ on the full ghost module.

\section{The square as an operator identity}

\begin{lemma}[Contractions of the charge]\label{lem:vir-Q-contractions}
The operator~\eqref{eq:vir-Q} satisfies
\begin{equation}\label{eq:vir-Q-contractions}
 \{Q,b_n\}=\mathcal L_n,\qquad \{Q,c_n\}=\nu_n.
\end{equation}
Moreover,
\begin{equation}\label{eq:vir-nu-identities}
 [\mathcal L_m,\nu_n]=(-2m-n)\nu_{m+n},\qquad
 [Q,\nu_n]=0.
\end{equation}
\end{lemma}

\begin{proof}
Contracting $b_n$ with the two $c$ factors in $R$ gives
\[
 \{R,b_n\}=\sum_{m\in\mathbb Z}(m-n):c_{-m}b_{m+n}:
 =\sum_{s\in\mathbb Z}(n+s):b_{n-s}c_s:
 =L_n^{\mathrm{gh}}.
\]
Contracting $c_n$ with the $b$ factor gives
\[
 \{R,c_n\}=-\frac12\sum_{r+s=-n}(r-s)c_{-r}c_{-s}=\nu_n.
\]
The matter term contributes $L_n^M$ to the first identity and zero to the second.
All contractions retain the same normal ordering of the uncontracted factors.

The commutator of $\mathcal L_m$ with $\nu_n$ uses~\eqref{eq:vir-ghost-transform}.
After reindexing, the coefficient of $c_uc_v$, $u+v=m+n$, is
\begin{align*}
 \frac12\bigl((u-v-m)(-m-u)+(u-v+m)(-m-v)\bigr)
 &=-\frac12(2m+n)(u-v).
\end{align*}
This proves the first part of~\eqref{eq:vir-nu-identities}.

For the second part, extend $\delta(c_n)=\nu_n$ as an odd derivation on products of $c$ modes.
Their mutual anticommutators are zero.
The coefficient of each alternating cubic monomial in $\delta^2(c_n)$ is a multiple of
\[
 (a-b)(a+b-d)+(b-d)(b+d-a)+(d-a)(d+a-b)=0.
\]
Thus $\delta^2(c_n)=0$.
The second identity in~\eqref{eq:vir-Q-contractions} identifies $[Q,\nu_n]$ with this expression.
Lemma~\ref{lem:vir-domain} permits these expansions on each algebraic input.
\end{proof}

\begin{theorem}[The Virasoro square]\label{thm:vir-square}
For every restricted matter module $M$ as above,
\begin{equation}\label{eq:vir-square}
 \boxed{\displaystyle
 Q^2=\frac{C-26}{12}\sum_{n>0}(n^3-n)c_{-n}c_n
 \quad\text{on }M\otimes\mathcal F.}
\end{equation}
In addition,
\begin{equation}\label{eq:vir-Q-L}
 [Q,\mathcal L_n]=-A(n)c_n,\qquad [Q,\mathcal L_0]=0.
\end{equation}
If $M\neq0$, the operator $Q$ is square-zero on the full state space if and only if $C=26$.
\end{theorem}

\begin{proof}
Put $D_m=[Q,\mathcal L_m]$.
The graded Jacobi identity,~\eqref{eq:vir-Q-contractions}, and~\eqref{eq:vir-total-central} give
\begin{align*}
 \{D_m,b_n\}
 &=[Q,[\mathcal L_m,b_n]]_{\mathrm s}
   -[\mathcal L_m,\{Q,b_n\}]\\
 &=(m-n)\mathcal L_{m+n}-[\mathcal L_m,\mathcal L_n]
 =-A(m)\delta_{m+n,0},\\
 \{D_m,c_n\}
 &=(-2m-n)\nu_{m+n}-[\mathcal L_m,\nu_n]=0.
\end{align*}
Hence $D_m+A(m)c_m$ supercommutes with all ghosts.
It has ghost degree one, so Lemma~\ref{lem:vir-commutant} makes it zero.
This proves~\eqref{eq:vir-Q-L}.

Let $K$ be the right side of~\eqref{eq:vir-square}.
A direct Clifford contraction gives
\[
 [K,b_n]=-A(n)c_n,\qquad [K,c_n]=0.
\]
For instance,
$[c_{-m}c_m,b_n]=\delta_{m+n,0}c_{-m}-\delta_{n-m,0}c_m$.
The oddness of $A$ gives the formula for both signs of $n$.
Also,
\[
 [Q^2,b_n]=[Q,\{Q,b_n\}]=[Q,\mathcal L_n],\qquad
 [Q^2,c_n]=[Q,\nu_n]=0.
\]
Thus $Q^2-K$ commutes with all ghosts and has ghost degree two.
Lemma~\ref{lem:vir-commutant} proves that it is zero.

When $C=26$, the displayed square vanishes.
For nonzero $v\in M$, the converse follows from
\[
 Q^2(v\otimes b_{-2}\Omega)
 =\frac{C-26}{2}\,v\otimes c_{-2}\Omega.
\]
The output is nonzero when $C\neq26$.
\end{proof}

\section{Fields, zero modes, and the intercept}

Fix a formal coordinate $z$ and define
\[
 b(z)=\sum_n b_nz^{-n-2},\qquad
 c(z)=\sum_n c_nz^{1-n},\qquad
 T^M(z)=\sum_nL_n^Mz^{-n-2}.
\]
The residue extracts the coefficient of $z^{-1}dz$.
For convergent Laurent series it equals the positively oriented integral divided by $2\pi i$.
All uses here are formal.
The two ghost contractions are the expansion of $(z-w)^{-1}$ in nonnegative powers of $w/z$.
For example,
\[
 \sum_{n\geq-1}z^{-n-2}w^{n+1}=\frac1{z-w}.
\]
The ghost stress tensor is
\begin{equation}\label{eq:vir-stress-field}
 T^{\mathrm{gh}}(z)=-2:b(z)\partial c(z):-:\partial b(z)c(z):
 =\sum_mL_m^{\mathrm{gh}}z^{-m-2}.
\end{equation}
Indeed, the coefficient of $b_rc_s$ is $r+2s$.
When $r+s=m$, this is $m+s$, as in~\eqref{eq:vir-ghost-L}.
Equation~\eqref{eq:vir-ghost-transform} specifies the conformal weights $2$ and $-1$.

Multiplying the fields and taking the residue gives exactly
\begin{equation}\label{eq:vir-charge-residue}
 Q=\operatorname{Res}_{z=0}
 \bigl(c(z)T^M(z)+:b(z)c(z)\partial c(z):\bigr)\,dz.
\end{equation}
For the matter term, the residue condition is $m+n=0$.
For the ghost term it is $r+s+t=0$.
Antisymmetrizing the two $c$ factors changes the coefficient $1-t$ to $(s-t)/2$, which gives~\eqref{eq:vir-ghost-R}.
If $\mathcal N$ means full normal ordering of the elementary ghost factors, then
\[
 \tfrac12\mathcal N(cT^{\mathrm{gh}})=:bc\partial c:.
\]
Thus~\eqref{eq:vir-charge-residue} also gives the residue of
$cT^M+\tfrac12\mathcal N(cT^{\mathrm{gh}})$.
Adding $\tfrac32\partial^2c$ to this current leaves its residue unchanged.
This specifies a zero mode, rather than a sum of all current modes.

\begin{proposition}[Intercept and vacuum conventions]\label{prop:vir-intercept}
For $t\in\mathbb C$, set $Q_t=Q+tc_0$.
Then
\begin{align}
 \{Q_t,b_n\}&=\mathcal L_n+t\delta_{n,0},
 & [Q_t,\mathcal L_0]&=0,
 \label{eq:vir-intercept-zero}\\
 Q_t^2&=\sum_{n>0}
 \left(\frac{C-26}{12}(n^3-n)-2tn\right)c_{-n}c_n.
 \label{eq:vir-intercept-square}
\end{align}
For $M\neq0$, absolute nilpotence holds exactly when $C=26$ and $t=0$.

Let $\downarrow=c_1\Omega$ and normal-order relative to this vacuum.
Its annihilators are $b_n$ for $n\geq0$ and $c_n$ for $n\geq1$.
Use a subscript $\downarrow$ for this ordering.
Then
\begin{equation}\label{eq:vir-intercept-dictionary}
 L_{n,\downarrow}^{\mathrm{gh}}=L_n^{\mathrm{gh}}+\delta_{n,0},
 \quad R_{\downarrow}=R+c_0,
 \quad
 \sum_n c_{-n}L_n^M+R_{\downarrow}-ac_0=Q+(1-a)c_0.
\end{equation}
In this convention the nilpotent intercept is $a=1$.
The underlying Clifford module is the same.
The vacuum $\downarrow$ has conformal energy $-1$ and ghost degree one in the $\Omega$ convention.
\end{proposition}

\begin{proof}
The first line follows from~\eqref{eq:vir-Q-contractions} and $[\mathcal L_0,c_0]=0$.
Since $c_0^2=0$,
\[
 Q_t^2=Q^2+t\{Q,c_0\},\qquad
 \{Q,c_0\}=\nu_0=-2\sum_{n>0}n c_{-n}c_n.
\]
This proves~\eqref{eq:vir-intercept-square}.
On $v\otimes\Omega$, its value is $-2t\,v\otimes c_{-1}c_1\Omega$.
Thus absolute nilpotence requires $t=0$.
Theorem~\ref{thm:vir-square} then requires $C=26$.

The change of vacuum reverses only the creator/annihilator pair $(c_1,b_{-1})$.
Its normal products satisfy
\[
 :b_{-1}c_1:_{\downarrow}-:b_{-1}c_1:
 =b_{-1}c_1+c_1b_{-1}=1.
\]
The coefficient of this pair in $L_0^{\mathrm{gh}}$ is $1$, and it occurs in no other $L_n^{\mathrm{gh}}$.
This gives the first identity in~\eqref{eq:vir-intercept-dictionary}.
The same contraction proof as in Lemma~\ref{lem:vir-Q-contractions} gives
\[
 \{R_{\downarrow}-R,b_n\}=\delta_{n,0},\qquad
 \{R_{\downarrow}-R,c_n\}=0.
\]
The difference $R_{\downarrow}-R-c_0$ vanishes by Lemma~\ref{lem:vir-commutant}.
The last identity in~\eqref{eq:vir-intercept-dictionary} follows.
The energy and degree of $\downarrow$ follow from those of $c_1$.
\end{proof}

\begin{proposition}[The nonnegative ghost polarization]\label{prop:vir-polarization-map}
Set $c^m=c_{-m}$, so that $\{b_m,c^n\}=\delta_m^n$.
Let $\mathcal F_{\mathrm{pol}}$ have vacuum $\Omega_{\mathrm{pol}}$ killed by $c^m$ for $m\leq0$ and by $b_m$ for $m>0$.
Its creators are $c^m$ for $m>0$ and $b_m$ for $m\leq0$.
Give its vacuum degree and energy zero, and assign these creators degrees $1,-1$ and energies $m,-m$ respectively.
There is a Clifford-module isomorphism
\[
 I:\mathcal F_{\mathrm{pol}}\longrightarrow\mathcal F,
 \qquad I(x\Omega_{\mathrm{pol}})=x\,c_0c_1\Omega,
\]
where $x$ is a product of the indicated creators.
It raises ghost degree by $2$ and lowers energy by $1$.
Let $R_{\mathrm{pol}}$ denote the cubic expression~\eqref{eq:vir-ghost-R} with this ordering.
Under $1_M\otimes I$,
\[
 \sum_m c^mL_m^M+R_{\mathrm{pol}}-ac^0
 \quad\longmapsto\quad Q+(1-a)c_0.
\]
Thus its critical charge term is $-c^0$, with both degree and energy shifts specified by $I$.
\end{proposition}

\begin{proof}
The vector $\Xi=c_0c_1\Omega$ obeys the required annihilator conditions.
For $m\leq0$, $c_{-m}\Xi=0$: $c_0$ and $c_1$ square to zero, and all $c_n$ with $n\geq2$ annihilate $\Xi$.
For $m>0$, $b_m\Xi=0$ by the Clifford relations and~\eqref{eq:vir-ghost-vacuum}.
Hence the displayed map intertwines the Clifford actions.
Away from the pairs $(b_0,c_0)$ and $(b_{-1},c_1)$ it identifies the same exterior creators.
On each of these two pairs, exterior multiplication and exterior differentiation exchange the empty and occupied states, with nonzero signs.
Their four-state tensor product is invertible.
This proves that $I$ is an isomorphism.
Since $\Xi$ has degree $2$ and energy $-1$, the grading statements follow.

Relative to $\Omega$, this ordering reverses both specified pairs.
The pair $(b_{-1},c_1)$ contributes $1$ to $L_0^{\mathrm{gh}}$, as in Proposition~\ref{prop:vir-intercept}.
The pair $(b_0,c_0)$ has coefficient zero there.
Consequently the transported ghost stress modes are $L_n^{\mathrm{gh}}+\delta_{n,0}$.
The contraction identities and Lemma~\ref{lem:vir-commutant} then give
$I R_{\mathrm{pol}}I^{-1}=R+c_0$.
The matter term commutes with $I$, which proves the charge formula.
\end{proof}

\section{Absolute and relative complexes}

Here absolute and relative refer to the $b_0$ constraint.
The central character of the Virasoro module is fixed, and the ghost variables are indexed by the modes $L_n$.
There is no ghost for the central generator in this carrier.

Set
\[
 \mathcal D_t^{\mathrm{rel}}=
 \ker(b_0:\mathcal D\to\mathcal D)
 \cap\ker(\mathcal L_0+t:\mathcal D\to\mathcal D).
\]
It inherits the ghost grading of $\mathcal D$.
The first identity in~\eqref{eq:vir-intercept-zero} and the commutation with $\mathcal L_0$ show that $Q_t$ preserves this subspace.
On it, the square is the restriction of~\eqref{eq:vir-intercept-square}.

\begin{corollary}[The critical complexes]\label{cor:vir-complexes}
At $C=26$ and $t=0$, the pairs
\[
 C_{\mathrm{abs}}^q(M)=M\otimes\mathcal F^q,
 \qquad
 C_{\mathrm{rel}}^q(M)=\mathcal D_0^{\mathrm{rel}}\cap(M\otimes\mathcal F^q)
\]
with differential $Q$ are cochain complexes of complex vector spaces.
The inclusion $C_{\mathrm{rel}}^\bullet(M)\to C_{\mathrm{abs}}^\bullet(M)$ is a chain map.
A homomorphism of restricted Virasoro modules induces chain maps by tensoring with $1_{\mathcal F}$.
If $\mathcal L_0$ acts semisimply, every nonzero weight $\lambda$ subcomplex of the absolute complex is contractible by $\lambda^{-1}b_0$.
\end{corollary}

\begin{proof}
The domain lemma defines the differential in every degree.
Theorem~\ref{thm:vir-square} gives its square, and~\eqref{eq:vir-intercept-zero} gives stability of the relative subspace.
A module homomorphism commutes termwise with~\eqref{eq:vir-Q}, hence with the finite sums on each state.
On the weight-$\lambda$ subspace,
$Q(\lambda^{-1}b_0)+(\lambda^{-1}b_0)Q=1$ by~\eqref{eq:vir-Q-contractions}.
\end{proof}

Relative nilpotence depends on the actual relative carrier.
The full-state converse in Theorem~\ref{thm:vir-square} cannot be restricted without checking its detecting vectors.

\begin{proposition}[Nonnegative matter weights]\label{prop:vir-relative-three-term}
Suppose $M=\bigoplus_{j\geq0}M_j$, with $L_0^M|_{M_j}=j$ and $L_n^MM_j\subseteq M_{j-n}$.
At $t=0$, for every central charge $C$, the relative carrier is
\begin{equation}\label{eq:vir-relative-carrier}
 \mathcal D_0^{\mathrm{rel}}
 =(M_0\otimes\mathbb C\Omega)
 \oplus(M_1\otimes\mathbb Cc_1\Omega)
 \oplus(M_0\otimes\mathbb Cc_{-1}c_1\Omega).
\end{equation}
Its differential is the three-term complex
\begin{equation}\label{eq:vir-relative-three-term}
 M_0\xrightarrow{\ L_{-1}^M\ } M_1
 \xrightarrow{\ L_1^M\ }M_0
\end{equation}
in ghost degrees $0,1,2$.
In particular, $Q^2=0$ on this relative carrier for every $C$.
\end{proposition}

\begin{proof}
The only negative ghost energy is that of $c_1$, and it can occur at most once.
Thus $\mathcal F$ has minimum energy $-1$.
Every $b$ creator has energy at least $2$, so no $b$ creator occurs at ghost energy $0$ or $-1$.
The complete bases at these two energies are
\[
 \mathcal F_{-1}=\langle c_1\Omega,c_0c_1\Omega\rangle,
 \qquad
 \mathcal F_0=\langle\Omega,c_0\Omega,c_{-1}c_1\Omega,c_{-1}c_0c_1\Omega\rangle.
\]
The condition $b_0=0$ removes precisely the terms containing $c_0$.
The total weight-zero condition therefore gives~\eqref{eq:vir-relative-carrier}.

The matter term and ghost term of~\eqref{eq:vir-Q} give
\begin{align*}
 Q(v\otimes\Omega)&=L_{-1}^Mv\otimes c_1\Omega &&(v\in M_0),\\
 Q(w\otimes c_1\Omega)&=L_1^Mw\otimes c_{-1}c_1\Omega &&(w\in M_1),\\
 Q(v\otimes c_{-1}c_1\Omega)&=0 &&(v\in M_0).
\end{align*}
In the middle line, $c_0L_0^Mw\otimes c_1\Omega$ cancels $w\otimes Rc_1\Omega=-w\otimes c_0c_1\Omega$.
For the last line, $R(c_{-1}c_1\Omega)=0$ by the two contractions in Lemma~\ref{lem:vir-Q-contractions}.
Finally,
$L_1^ML_{-1}^Mv=2L_0^Mv+L_{-1}^ML_1^Mv=0$ on $M_0$.
This proves~\eqref{eq:vir-relative-three-term} and its nilpotence.
\end{proof}

For example, the universal Virasoro vacuum module of central charge $C$ is induced from a vector $\mathbf1$ killed by $L_n$ for $n\geq-1$.
The central generator acts by $C$.
The ordered monomials in $L_{-n}$, $n\geq2$, form its basis by the Poincar\'e--Birkhoff--Witt theorem.
Hence $M_0=\mathbb C\mathbf1$ and $M_1=0$.
Its relative complex has basis $\mathbf1\otimes\Omega,\mathbf1\otimes c_{-1}c_1\Omega$ and zero differential for every $C$.
Trivial matter gives the same relative complex at $C=0$.
Its absolute square is nonzero by~\eqref{eq:vir-first-square}.
At the other boundary, suppose a restricted module contains $v\neq0$ with $L_0^Mv=-2v$.
Then $v\otimes b_{-2}\Omega$ lies in $\mathcal D_0^{\mathrm{rel}}$ and
\[
 Q^2(v\otimes b_{-2}\Omega)=\frac{C-26}{2}v\otimes c_{-2}\Omega.
\]
This relative square is nonzero when $C\neq26$.
Thus the critical central charge guarantees relative nilpotence for every restricted module, while particular relative carriers can have further vanishing.

\section{Conformal vertex algebra states}

Let $V$ be an ordinary conformal vertex algebra over $\mathbb C$ with conformal vector $\omega$.
Write
$Y(\omega,z)=\sum_nL_n^Vz^{-n-2}$ and let $C$ be its central charge.
The lower truncation axiom makes $V$ a restricted matter module.
Consequently~\eqref{eq:vir-Q} defines the degree-one operator on the full algebraic state carrier $V\otimes\mathcal F$.
Equation~\eqref{eq:vir-charge-residue} identifies it with the zero mode of the specified normal-ordered current.
At $C=26$, Corollary~\ref{cor:vir-complexes} supplies both the absolute complex and its $b_0$-relative subcomplex.
When $V$ has nonnegative integral conformal weights, Proposition~\ref{prop:vir-relative-three-term} computes the relative complex directly.

The ghost degree here is the cohomological degree.
Conformal weight is the eigenvalue of $L_0^V+L_0^{\mathrm{gh}}$ and is independent of ghost degree.
The residue uses the fixed formal coordinate $z$.
A comparison with a complex of collision-state sheaves requires maps in that sheaf category, with compatible restrictions and differentials.
The state-space construction above supplies its Virasoro differential and its square.
